% ============================================================
% Chapter 4 — Burp Suite Interface Explained
% ============================================================
\chapter{Burp Suite Interface Explained}
\label{ch:burpinterface}

\noteblock{This chapter explains every Burp Suite module you will use in this course. Read this before starting any attack module — it is your reference for understanding what each tool does and when to use it.}

% ─────────────────────────────────────────────────────────
\section{Proxy}
\label{sec:proxy}

\subsection{What It Does}

The Proxy is the heart of Burp Suite. It intercepts all HTTP/HTTPS traffic between your browser and the target application, allowing you to read, pause, modify, and forward every request and response.

\subsection{Traffic Flow}

\begin{center}
\begin{tabular}{ccccc}
\textbf{Browser} & $\longrightarrow$ & \textbf{Burp Proxy} & $\longrightarrow$ & \textbf{Target App} \\
& & \updownarrow & & \\
& & \textbf{You (tester)} & & \\
\end{tabular}
\end{center}

\subsection{Intercept ON / OFF}

\begin{table}[H]
\centering
\caption{Intercept mode comparison}
\begin{tabularx}{\textwidth}{>{\bfseries}lX}
\toprule
Mode & Behavior \\
\midrule
Intercept ON & Each request pauses in Burp — you must manually forward or drop it \\
Intercept OFF & Traffic flows through transparently — all requests logged in HTTP History \\
\bottomrule
\end{tabularx}
\end{table}

\noteblock{For passive reconnaissance (browsing to map endpoints), use \textbf{Intercept OFF}. Turn it \textbf{ON} only when you need to modify a specific request before it reaches the server.}

\subsection{HTTP History}

HTTP History (\textbf{Proxy} $\to$ \textbf{HTTP History}) records every request/response passing through Burp. Key columns:

\begin{table}[H]
\centering
\caption{HTTP History columns explained}
\begin{tabular}{ll}
\toprule
Column & Meaning \\
\midrule
\# & Request number (sequential) \\
Method & GET, POST, PUT, DELETE \\
URL & Full request URL \\
Status & HTTP response code (200, 404, 302, etc.) \\
Length & Response body size (useful for spotting differences) \\
MIME & Content type of the response \\
\bottomrule
\end{tabular}
\end{table}

\subsection{Request Modification}

When Intercept is ON, you can modify the request before forwarding:
\begin{itemize}
  \item Change parameter values
  \item Add, remove, or modify headers
  \item Change the HTTP method (GET $\to$ POST)
  \item Inject payloads into the body
\end{itemize}

\subsection{Response Analysis}

Click any request in HTTP History $\to$ click \textbf{Response} tab to see:
\begin{itemize}
  \item HTTP status code and reason phrase
  \item All response headers (look for missing security headers)
  \item Response body (HTML, JSON, XML)
  \item Use \textbf{Ctrl+F} to search for keywords: \texttt{password}, \texttt{token}, \texttt{admin}, \texttt{error}
\end{itemize}

\screenshotbox{Burp Proxy — HTTP History with request/response panes}

% ─────────────────────────────────────────────────────────
\section{Target}
\label{sec:target}

\subsection{What It Does}

The Target module maintains a \textbf{Site Map} of everything discovered about the target application and lets you control the \textbf{Scope} — what Burp tracks and tests.

\subsection{Setting Scope}

Always define scope before testing. This prevents accidentally testing out-of-scope systems.

\begin{enumerate}
  \item Go to \textbf{Target} $\to$ \textbf{Scope settings}
  \item Click \textbf{Add}
  \item Enter: \texttt{https://recon.lab.uzyntra.com} (or your target)
  \item Optionally: set \textbf{Out-of-scope} entries for paths you must avoid
\end{enumerate}

\warnblock{In a real engagement, only test what is explicitly in scope. Testing out-of-scope systems — even accidentally — can terminate an engagement and create legal issues.}

\subsection{Site Map}

The Site Map (\textbf{Target} $\to$ \textbf{Site Map}) builds automatically as you browse. It shows:
\begin{itemize}
  \item All discovered URLs and endpoints
  \item HTTP methods seen at each endpoint
  \item Response codes
  \item Linked resources (JS files, CSS, images)
\end{itemize}

\exerciseblock{
\textbf{Exercise 4.1 — Target Mapping}\\[4pt]
\begin{enumerate}
  \item Add \texttt{https://recon.lab.uzyntra.com} to Burp scope
  \item Turn Intercept \textbf{OFF}
  \item Browse: \texttt{/}, \texttt{/about}, \texttt{/contact}, \texttt{/debug}, \texttt{/admin}, \texttt{/backup}, \texttt{/robots.txt}
  \item Open Target $\to$ Site Map
  \item Expand the host — how many endpoints were discovered?
  \item Right-click the host $\to$ \textbf{Spider this host} to find more
\end{enumerate}
}

% ─────────────────────────────────────────────────────────
\section{Repeater}
\label{sec:repeater}

\subsection{Purpose}

Repeater allows you to send a captured HTTP request multiple times with modifications, comparing the responses. It is the most used Burp tool for manual vulnerability testing.

\subsection{Sending a Request to Repeater}

\begin{enumerate}
  \item In HTTP History or Intercept, right-click a request
  \item Select \textbf{Send to Repeater} (or press \textbf{Ctrl+R})
  \item Switch to the \textbf{Repeater} tab
  \item Modify the request in the left pane
  \item Click \textbf{Send}
  \item Analyze the response in the right pane
\end{enumerate}

\screenshotbox{Repeater — request on left, response on right}

\subsection{Common Repeater Actions}

\begin{table}[H]
\centering
\caption{Repeater testing steps}
\begin{tabularx}{\textwidth}{>{\bfseries}lXX}
\toprule
Step & Action & Security Meaning \\
\midrule
Change parameter value & Edit \texttt{id=1} to \texttt{id=2} & Test for IDOR \\
Remove Authorization header & Delete the \texttt{Authorization:} line & Test for missing auth \\
Change HTTP method & GET $\to$ POST & Test for method confusion \\
Inject SQL & Change value to \texttt{' OR 1=1 --} & Test for SQL injection \\
Modify Content-Type & Change \texttt{application/json} to \texttt{text/xml} & Test for parser confusion \\
\bottomrule
\end{tabularx}
\end{table}

% ─────────────────────────────────────────────────────────
\section{Intruder}
\label{sec:intruder}

\subsection{Purpose}

Intruder automates customized attacks — iterating through payloads in request parameters. Use cases: brute force, fuzzing, enumeration.

\subsection{Attack Types}

\begin{table}[H]
\centering
\caption{Burp Intruder attack types}
\begin{tabularx}{\textwidth}{>{\bfseries}lX}
\toprule
Type & Description \\
\midrule
Sniper & One payload position, one wordlist — most common \\
Battering Ram & Same payload inserted into all positions simultaneously \\
Pitchfork & Multiple positions, multiple lists — pairs items (e.g. user+pass) \\
Cluster Bomb & Multiple positions, all combinations — combinatorial (slow) \\
\bottomrule
\end{tabularx}
\end{table}

\subsection{Setting Positions}

\begin{enumerate}
  \item Send a request to Intruder (\textbf{Ctrl+I})
  \item Go to \textbf{Positions} tab
  \item Click \textbf{Clear §} to remove all auto-detected positions
  \item Highlight the value you want to fuzz (e.g. a username or directory name)
  \item Click \textbf{Add §} — the value becomes \texttt{\S{}payload\S{}}
\end{enumerate}

\subsection{Community Edition Limitation}

\dangerblock{Burp Suite Community Edition \textbf{throttles Intruder to 1 request per second}. For large wordlists (10,000+ entries), this is very slow. For real engagements, use Professional Edition or alternative tools like \texttt{ffuf} or \texttt{hydra} for brute forcing.}

% ─────────────────────────────────────────────────────────
\section{Decoder}
\label{sec:decoder}

\subsection{Purpose}

Decoder encodes and decodes data in various formats. Essential for analyzing tokens, cookies, and obfuscated values.

\subsection{Supported Formats}

\begin{table}[H]
\centering
\caption{Decoder supported encodings}
\begin{tabularx}{\textwidth}{>{\bfseries}lXl}
\toprule
Format & Use Case & Example \\
\midrule
URL Encode & Decode URL parameters & \texttt{\%3D} $\to$ \texttt{=} \\
Base64 & Decode JWT headers/payloads, cookies & \texttt{dXNlcg==} $\to$ \texttt{user} \\
HTML Entity & Decode XSS-filtered output & \texttt{\&lt;} $\to$ \texttt{<} \\
Hex & Decode binary or hex-encoded data & \texttt{48656c6c6f} $\to$ \texttt{Hello} \\
Gzip & Decompress compressed responses & binary $\to$ text \\
\bottomrule
\end{tabularx}
\end{table}

\subsection{Using Decoder}

\begin{enumerate}
  \item Go to \textbf{Decoder} tab
  \item Paste the encoded value
  \item Select \textbf{Decode as} $\to$ choose format
  \item View the decoded output
  \item Stack operations: decode Base64 $\to$ then URL decode the result
\end{enumerate}

\noteblock{JWT tokens use Base64URL encoding (slightly different from standard Base64 — uses \texttt{-} and \texttt{\_} instead of \texttt{+} and \texttt{/}). Burp's Decoder handles both variants.}

% ─────────────────────────────────────────────────────────
\section{Comparer}
\label{sec:comparer}

\subsection{Purpose}

Comparer performs a word-level or byte-level diff between two responses. This is invaluable when you need to spot subtle differences — for example, determining whether a username exists (different response length or content).

\subsection{Using Comparer}

\begin{enumerate}
  \item In HTTP History, right-click a response $\to$ \textbf{Send to Comparer}
  \item Repeat with a second response
  \item Go to \textbf{Comparer} tab
  \item Click \textbf{Words} for semantic diff or \textbf{Bytes} for binary diff
  \item Highlighted differences show what changed between the two responses
\end{enumerate}

\exerciseblock{
\textbf{Exercise 4.2 — Comparer Practice}\\[4pt]
\begin{enumerate}
  \item In Repeater, send a request to \texttt{/admin} with no cookies
  \item Send it to Comparer
  \item Add an \texttt{Authorization: Bearer fake\_token} header and send again
  \item Send that response to Comparer
  \item Compare — what changed in the response?
\end{enumerate}
}

% ─────────────────────────────────────────────────────────
\section{Dashboard}

The Dashboard provides:
\begin{itemize}
  \item \textbf{Event log} — all Burp activities and errors
  \item \textbf{Issue activity} — vulnerability findings (Pro only for auto-scan)
  \item \textbf{Tasks} — manage running operations
\end{itemize}

\noteblock{In Community Edition, the Dashboard mostly shows the Event log. Check it if Burp behaves unexpectedly — connection errors, certificate issues, and configuration problems appear here.}
